\documentclass[aspectratio=169]{beamer}
\usepackage[numbers]{natbib}
\usepackage{booktabs}

\usetheme{default}
\setbeamertemplate{navigation symbols}{}
\setbeamertemplate{footline}[frame number]
\setbeamercolor{frametitle}{fg=black}
\setbeamerfont{frametitle}{series=\bfseries}
\setbeamertemplate{itemize items}[circle]

% natbib inside beamer: keep citations unobtrusive
\setbeamercolor{bibliography entry author}{fg=black}
\renewcommand{\bibsection}{}

\title{Where LLM Agents Break}
\subtitle{A literature review of the reason--act loop and its failure modes}
\author{}
\date{}

\begin{document}

\begin{frame}[plain]
  \titlepage
\end{frame}

\begin{frame}{The claim of this review}
  \begin{itemize}
    \item Three papers define the modern LLM agent: an architecture, a
      self-improvement mechanism, and an evaluation.
    \item Read together, they do \emph{not} describe steady progress.
    \item They describe a paradigm whose two central failures have no
      principled fix, and a metric showing the real bottleneck is
      \textbf{consistency}, not capability.
    \item[]
    \item \textbf{Thesis:} the field improves peak ability while the binding
      constraint is variance, and no method recovers a trajectory once it
      derails.
  \end{itemize}
\end{frame}

\begin{frame}{The architecture: interleave thought and action}
  ReAct \citep{yao2023reactsynergizingreasoningacting} augments the action space
  to $\hat{A} = A \cup L$. A \emph{thought} is an action in language space that
  ``does not affect the external environment, thus leading to no observation
  feedback'' --- it only updates context.

  \vspace{1em}
  Motivation: chain-of-thought is ``a static black box, in that the model uses
  its own internal representations to generate thoughts and is not grounded in
  the external world.''

  \vspace{1em}
  Frozen PaLM-540B, 1--6 in-context examples. No training.
\end{frame}

\begin{frame}{Grounding works --- and costs something}
  \begin{center}
  \begin{tabular}{lrr}
    \toprule
    Judged failure mode & ReAct & CoT \\
    \midrule
    Hallucination        & 0\%  & 56\% \\
    False positive       & 6\%  & 14\% \\
    Reasoning error      & 47\% & 16\% \\
    Search result error  & 23\% & --- \\
    \bottomrule
  \end{tabular}
  \end{center}

  \vspace{0.8em}
  Hallucination is eliminated. Reasoning errors nearly triple, because ``such a
  structural constraint also reduces its flexibility in formulating reasoning
  steps.''

  \vspace{0.5em}
  Net effect on HotpotQA: ReAct \textbf{27.4} vs CoT \textbf{29.4} exact match
  --- the paradigm \emph{loses}. It wins on FEVER, 60.9 vs 56.3.
\end{frame}

\begin{frame}{The two failures with no principled fix}
  \begin{itemize}
    \item \textbf{Loop entrapment.} ``the model repetitively generates the
      previous thoughts and actions'' and cannot break out. Folded into the 47\%
      reasoning-error bucket, so its own rate was never reported. Attributed to
      greedy decoding --- untested.
    \item \textbf{Unrecoverable retrieval failure.} 23\% of errors. The damage is
      not bad retrieval but the absence of any backtracking mechanism.
  \end{itemize}

  \vspace{1em}
  Both are made terminal by a hard step budget (7 steps HotpotQA, 5 FEVER) plus
  the finding that ``more steps will not improve ReAct performance.''

  \vspace{0.5em}
  \textbf{More compute does not help. Only trajectory recovery would.}
\end{frame}

\begin{frame}{The mechanism: learn without weights}
  Reflexion \citep{shinn2023reflexionlanguageagentsverbal} reinforces agents
  ``not by updating weights, but instead through linguistic feedback.'' A failed
  trial becomes a verbal summary in episodic memory, prepended to the next
  attempt.

  \vspace{1em}
  Results are strong: \textbf{130 of 134} ALFWorld tasks; \textbf{91\%} pass@1 on
  HumanEval against a reported \textbf{80\%} for GPT-4.

  \vspace{1em}
  Baseline ReAct plateaus --- performance ``halts between trials 6 and 7,''
  converging at a 22\% hallucination rate ``with no signs of long-term
  recovery.''
\end{frame}

\begin{frame}{But the mechanism is thinner than the framing}
  \begin{itemize}
    \item \textbf{The trigger is an \texttt{if} statement.} Reflection fires when
      the agent repeats an action ``more than 3 cycles,'' or ``exceeds 30''
      actions. Two unablated constants, tuned to one domain --- and this is a
      \emph{loop detector}, i.e.\ a hardcoded patch for ReAct's signature bug.
    \item \textbf{``Long-term memory'' holds three entries.} $\Omega$ is
      ``usually set to 1-3,'' truncated to ``the last 3 self-reflections'' to fit
      the context window --- yet learning is reported over 12 trials.
    \item \textbf{Reflection can entrench a wrong diagnosis.} Acknowledged: it
      ``may still succumb to non-optimal local minima solutions.'' Nothing
      detects a bad reflection.
  \end{itemize}
\end{frame}

\begin{frame}{The unstated scope limit}
  The loop requires ``reset the environment, and start a new trial.''

  \vspace{1em}
  Every reported gain depends on retrying from a clean state. That silently
  excludes:
  \begin{itemize}
    \item irreversible actions --- a sent message, a payment, a physical movement
    \item tasks with no machine-checkable success signal
  \end{itemize}

  \vspace{1em}
  Neither appears in the paper's Limitations section. Together they rule out most
  settings agents are actually being deployed into.
\end{frame}

\begin{frame}{The evaluation: measure consistency, not peak}
  $\tau$-bench \citep{yao2024taubenchbenchmarktoolagentuserinteraction} proposes
  \textbf{pass$^k$} --- ``the chance that all $k$ i.i.d.\ task trials are
  successful.'' The deliberate inverse of pass@$k$, ``the chance that at least
  one out of $k$ \dots\ is successful.''

  \vspace{1em}
  \begin{center}
  \begin{tabular}{lr}
    \toprule
    gpt-4o, function calling & success \\
    \midrule
    pass$^1$ & $< 50\%$ \\
    pass$^8$ (retail) & $< 25\%$ \\
    \bottomrule
  \end{tabular}
  \end{center}

  \vspace{0.8em}
  An agent that can usually do a task often cannot do it \emph{eight times
  running}. Averaged single-run benchmarks cannot see this at all.
\end{frame}

\begin{frame}{And the benchmark grades itself generously}
  Reward is terminal-state only: the final database must be ``identical to the
  unique ground truth outcome database.''

  \vspace{1em}
  The authors state the consequence plainly:

  \begin{quote}
    ``$r = 1$ might be a necessary but not sufficient condition for a successful
    episode e.g., the agent might issue the return without explicit user
    confirmation, which violates the policy.''
  \end{quote}

  \vspace{0.5em}
  A policy-violating trajectory can score 1. So $< 50\%$ is an \textbf{upper
  bound} on compliance --- the true numbers are worse.
\end{frame}

\begin{frame}{Recurring failures, not per-paper flaws}
  \begin{itemize}
    \item \textbf{Load-bearing magic numbers.} ReAct: backoff at 7/5 steps, vote
      threshold $n/2$. Reflexion: 3 cycles, 30 actions. All hand-set, none
      ablated.
    \item \textbf{Nothing recovers within a trajectory.} Two options exist:
      continue and fail, or discard the episode and restart. Mid-trajectory
      repair appears in none of the three.
    \item \textbf{Measurement flatters the method.} ReAct: 29\% label ambiguity
      against a 2-point gap. Reflexion: 91\% vs 80\% compares multi-trial to
      single-shot. $\tau$-bench: reward passes policy violations.
    \item \textbf{Retry/oracle assumptions exclude deployment.} Methods are built
      where retries are free; reliability is measured where they are not.
  \end{itemize}
\end{frame}

\begin{frame}{Where the openings are}
  \begin{enumerate}
    \item \textbf{Trajectory-aware reward for $\tau$-bench.} Check policy
      compliance along the path, not just the terminal database state. Motivated
      by the authors' own admission; would tighten every number in the paper.
      \emph{Best-scoped entry point.}
    \item \textbf{Self-assessed reflection triggers.} Does Reflexion's gain
      survive when the agent decides for itself when to reflect? Never tested.
    \item \textbf{Mid-trajectory recovery.} Backtracking over the
      thought--action trace. The largest structural gap, and the most likely to
      generalize.
    \item \textbf{Separate agent variance from user-simulator variance.}
      $\tau$-bench does not; without that split, it is unclear how much of the
      pass$^8$ collapse is even the agent's fault.
  \end{enumerate}
\end{frame}

\begin{frame}[allowframebreaks]{References}
  \bibliographystyle{plainnat}
  \bibliography{../../refs}
\end{frame}

\end{document}
